%% sec-9-conclusion.tex -- Conclusion

\section{Conclusion}
\label{sec:conclusion}

We have systematised more than 180 representative works on
real-time dynamic 4DGS rendering on mobile GPUs published between
2023 and the first half of 2026. The survey is organised around
four themes: representation, training-time acceleration,
rendering-time acceleration, and mobile deployment. Current
Snapdragon~8~Elite devices already achieve desktop-class real-time
3DGS, but 4DGS remains limited by a triple bottleneck of memory
bandwidth, Vulkan driver maturity, and the mobile power wall.

Promising future directions include feed-forward 4DGS,
hardware--software co-design, 4DGS ray tracing, a temporal
compression standard, and cross-modal fusion. We expect the
emergence of a community-wide \emph{4DGS compression standard}---a
codec profile that combines dynamic/static separation, residual
coding, and rate--distortion optimisation; preliminary prototypes
along this line include 4DGS-CC~\cite{chen20254dgscc} and the
feed-forward densification policy of
L2D2-GS~\cite{song2026l2d2gs}, both of which could anchor such a
standard. This survey is the academic counterpart of a separate
engineering roadmap (the authors' production-track 4DGS
renderer); the survey emphasises breadth across the research
landscape rather than the engineering decisions of any one
project.

\subsection{Limitations of This Survey}
We acknowledge three known limitations of the present survey.
First, in \textbf{scope}: we cover 2023 to the first half of
2026; very early 4DGS works from 2022 (e.g.\ the original NeRF-t
lineage) and any submission that appears after the cut-off are
excluded. Second, in \textbf{methodology}: the survey is
English-only and prioritises code-available or reproducible
papers; non-English venues and theoretical-only submissions are
underrepresented. Third, in \textbf{known gaps}: (i)~Faction~A
(the earliest 4DGS canonical-space methods) receives less space
than the newer feed-forward and compression factions;
(ii)~feed-forward 4DGS is covered via L2D2-GS but the parallel
diffusion-based generation literature is only sampled; and
(iii)~cross-domain fusion papers (4DGS + SLAM,
4DGS + simulation) are sparse and we surface them only
in~\S\ref{sec:discussion}.

\subsection{Open Science and Reading Notes}
In the spirit of survey reproducibility, the authors maintain
detailed per-paper notes (not included in this submission) that
record the source PDF page for every quantitative claim, the
citation chain, and a faction-classified summary. Every
\texttt{\textbackslash cite} key in this PDF corresponds to a
paper note in this private collection, which the authors are
happy to share on request. Readers who spot missing papers,
mis-categorised factions, or updated performance numbers are
invited to contact the corresponding author.

\bigskip
\noindent\textbf{Section reading notes}:
the per-paper notes collection.
